% =========================================================================
\section{Seminarios prácticos}\label{sec:seminarios}
% =========================================================================
Esta sección resuelve, con todo detalle, los siete problemas de la relación del tema~\cite{PROBLEMAS-T3}, más un ejemplo modelo del método de los
coeficientes de potencial~\cite{GUIA-T3}. Cada problema sigue la misma estructura que en los temas anteriores: el \emph{enunciado}, la \emph{idea} que guía
la resolución, la \emph{resolución} paso a paso y un comentario sobre la \emph{interpretación} del resultado y sus comprobaciones. Salvo que se diga lo
contrario, $\phi(\infty)=0$ y $\kc=1/(4\pi\eo)$.

% -------------------------------------------------------------------------
\subsection{Cómo plantear un problema de conductores}\label{sem:metodo}
% -------------------------------------------------------------------------
Todo problema de conductores se reduce a un sistema lineal: \emph{cada conductor aporta una ecuación y una incógnita}, porque de cada uno se conoce o su
carga o su potencial. La primera tarea es, pues, traducir el enunciado (\cref{tab:traduccion}); la segunda, escoger el método (\cref{tab:mapa}).

\begin{table}[ht]
\centering
\caption{Traducción de la situación física a ecuaciones.}
\label{tab:traduccion}
\small
\begin{tabularx}{\linewidth}{@{}>{\raggedright\arraybackslash}p{4.1cm}>{\raggedright\arraybackslash}X>{\raggedright\arraybackslash}p{4.3cm}@{}}
  \toprule
  Situación & Dato & Incógnita\\
  \midrule
  Aislado & $Q_i$ conocida (se conserva) & $\phi_i$\\
  Conectado a tierra & $\phi_i=0$ & $Q_i$ (fluye carga desde o hacia tierra)\\
  Unido a $j$ por un hilo fino & $\phi_i=\phi_j$ y $Q_i+Q_j$ conservada & el reparto $Q_i$, $Q_j$\\
  Batería $V_0$ entre $i$ y $j$ & $\phi_i-\phi_j=V_0$ y $Q_i+Q_j$ conservada & $Q_i$ (y $Q_j=-Q_i$ si eran neutros)\\
  Se conecta y se desconecta & tras desconectar, la carga queda congelada & la carga de cada etapa\\
  \bottomrule
\end{tabularx}
\end{table}

\begin{table}[ht]
\centering
\caption{Mapa de decisión: qué método conviene según lo que pide el enunciado.}
\label{tab:mapa}
\small
\begin{tabularx}{\linewidth}{@{}>{\raggedright\arraybackslash}X>{\raggedright\arraybackslash}p{5.6cm}c@{}}
  \toprule
  Lo que pide el enunciado & Método & Problemas\\
  \midrule
  Potencial o campo de un conductor aislado; ruptura dieléctrica; efecto punta & M1. Conductor aislado & 3.1\\
  Esferas o cilindros concéntricos con cargas, tierras o baterías & M2. Cadena de Gauss y potenciales & 3.2, 3.4\\
  «Calcula los coeficientes de capacidad e influencia» & M3. Coeficientes (caminos A y B) & 3.3\\
  Placas planas paralelas; cargas en cada cara & M4. Superposición de planos & 3.5\\
  Esferas pequeñas y lejanas ($d\gg R$); tierras sucesivas & M5. Esferas lejanas con $p_{ij}$ aproximados & 3.6, 3.7\\
  \bottomrule
\end{tabularx}
\end{table}

\begin{receta}[Los cinco métodos]
\begin{enumerate}[label=\textbf{M\arabic*.}, leftmargin=2.2em]
  \item \emph{Conductor aislado.} $\phi=\kc Q/R$, $E_{\mathrm{sup}}=\sigma/\eo=\phi/R$ (esfera).
  \item \emph{Concéntricos.} (1) Cargas en cada cara con la regla de las caras (\cref{prop:caras}); (2) Gauss en cada hueco: $E=\kc\Qenc/r^{2}$ (o
        $\lambda_{\mathrm{enc}}/2\pi\eo r$), con $\Qenc$ la suma de las cargas netas interiores, y $\vE=\vc{0}$ dentro del metal; (3) potencial del más
        externo, $\phi_N=\kc Q_{\mathrm{total}}/R_N$; (4) cadena hacia dentro, $\phi_k-\phi_{k+1}=\kc\Qenc\bigl(\frac{1}{R_k}-\frac{1}{R_{k+1}}\bigr)$;
        (5) imponer tierras, baterías y aislamientos.
  \item \emph{Coeficientes.} Camino A (directo): $\phi_j=1$, los demás a tierra; las cargas obtenidas son la columna $j$ de $\Cm$. Camino B: $Q_j=1$,
        los demás descargados; los potenciales son la columna $j$ de $\Pm$; luego se invierte.
  \item \emph{Placas paralelas.} Cada cara es un plano que crea $\pm\sigma/2\eo$; se impone conservación de carga y $\vE=\vc{0}$ dentro de cada placa.
  \item \emph{Esferas lejanas.} $p_{ii}\simeq\kc/R_i$, $p_{ij}\simeq\kc/d_{ij}$; en procesos por etapas, solo cambia la carga de la esfera conectada.
\end{enumerate}
\end{receta}

% =========================================================================
\subsection{Problema 3.1. Globo conductor y campo de ruptura}\label{sem:globo}
% =========================================================================
\begin{problema}[Potencial máximo de un globo]\label{prob:globo}
Supongamos un globo esférico de radio $5\,\mathrm{cm}$ recubierto por una fina película conductora. Si el campo electrostático de ruptura en la atmósfera
es de $3\cdot10^{6}\,\mathrm{V/m}$, calcula el máximo potencial al que puede conectarse el globo.
\end{problema}

\begin{resolucion}
\paso{Idea.} La \emph{ruptura dieléctrica} es el fenómeno por el que un aislante (aquí el aire) se ioniza y se convierte en un medio conductor cuando el
campo supera un valor crítico $\Erup$: el globo empezaría a descargarse a través del aire. Como el campo es máximo justo en la superficie del conductor, la
condición de ruptura hay que imponerla allí. Es el método M1.

\paso{Paso 1: campo y potencial del globo.} La película conductora es una esfera conductora de radio $R$. Con carga $Q$, la carga se reparte
uniformemente y, por Gauss, el campo exterior es el de una carga puntual en el centro:
\[
  \vE(\vc{r})=\frac{1}{4\pi\eo}\,\frac{Q}{r^{2}}\,\ur\quad(r>R),\qquad E(R)=\frac{1}{4\pi\eo}\,\frac{Q}{R^{2}},\qquad
  \phi=\frac{1}{4\pi\eo}\,\frac{Q}{R} .
\]
\paso{Paso 2: relación entre campo superficial y potencial.} Dividiendo $E(R)$ entre $\phi$ desaparecen la carga y la constante:
\[
  \frac{E(R)}{\phi}=\frac{\kc Q/R^{2}}{\kc Q/R}=\frac{R}{R^{2}}=\frac1R\quad\Longrightarrow\quad\resalta{\phi=E(R)\,R .}
\]
\paso{Paso 3: condición de ruptura.} Como $E(r)$ decrece fuera del globo, el campo más intenso está en la superficie. Si $E(R)$ supera $\Erup$, el aire se
ioniza y el globo se descarga. Por tanto,
\[
  \phi_{\max}=\Erup\,R=3\cdot10^{6}\,\frac{\mathrm{V}}{\mathrm{m}}\cdot5\cdot10^{-2}\,\mathrm{m}=1{,}5\cdot10^{5}\,\mathrm{V}
  \qquad\Longrightarrow\qquad\boxed{\;\phi_{\max}=150\,\mathrm{kV}\;}
\]
\end{resolucion}

% =========================================================================
%  Figura (Problema 3.1): globo conductor en el umbral de ruptura.
%  (a) E/E_rup y phi/phi_max frente a r/R: E = (R/r)^2, phi = R/r fuera.
%  (b) phi_max = E_rup R frente a R (escala logarítmica), E_rup = 3e6 V/m.
% =========================================================================
\begin{figure}[t]
\centering
\begin{tikzpicture}[font=\small]
  \begin{axis}[tema, name=a, width=7.0cm, height=5.0cm, xmin=0, xmax=4.2, ymin=0, ymax=1.25,
      xlabel={$r/R$}, ylabel={magnitud normalizada}, xtick={0,1,2,3,4}, ytick={0,0.5,1},
      legend style={at={(0.97,0.72)}, anchor=north east}]
    \fill[azulvelo!60] (axis cs:0,0) rectangle (axis cs:1,1.25);
    \node[font=\scriptsize, text=azulprof, rotate=90] at (axis cs:0.5,0.62) {globo};
    \addplot[azulprof, line width=1.2pt, domain=1:4.2, samples=60] {1/x^2};
    \addlegendentry{$E/\Erup$}
    \addplot[azulacero, line width=1.2pt, densely dashed, domain=1:4.2, samples=60] {1/x};
    \addlegendentry{$\phi/\phi_{\max}$}
    \addplot[azulprof, line width=1.2pt, forget plot] coordinates {(0,0) (1,0)};
    \addplot[azulacero, line width=1.2pt, densely dashed, forget plot] coordinates {(0,1) (1,1)};
    \draw[cargapos, line width=0.7pt, dashdotted] (axis cs:0,1) -- (axis cs:4.2,1);
    \node[font=\scriptsize, text=cargapos, anchor=south east] at (axis cs:4.15,1.01) {umbral de ruptura};
  \end{axis}
  \node[font=\small, anchor=north] at ($(a.outer south)+(0,-0.05)$) {(a)};
  \begin{axis}[tema, name=b, at={($(a.east)+(1.6cm,0)$)}, anchor=west, width=7.0cm, height=5.0cm,
      xmode=log, ymode=log, xmin=1e-4, xmax=1, ymin=1e2, ymax=1e7,
      xlabel={$R$ (m)}, ylabel={$\phi_{\max}$ (V)}, log basis x=10, log basis y=10]
    \addplot[azulprof, line width=1.2pt, domain=1e-4:1, samples=20] {3e6*x};
    \fill[cargapos] (axis cs:0.05,1.5e5) circle (2pt);
    \node[font=\scriptsize, anchor=north west, text=cargapos] at (axis cs:0.055,1.5e5) {$150\,\mathrm{kV}$};
    \node[font=\scriptsize, anchor=north west, text=grispizarra] at (axis cs:1.3e-4,5e6) {$\phi_{\max}=\Erup R$};
  \end{axis}
  \node[font=\small, anchor=north] at ($(b.outer south)+(0,-0.05)$) {(b)};
\end{tikzpicture}
\caption{Problema 3.1. (a) Campo y potencial de una esfera conductora cargada justo en el umbral de ruptura: el campo es máximo en la superficie y decae como
$1/r^{2}$, de modo que la ruptura empieza allí. (b) Potencial máximo frente al radio: cuanto más pequeño es el conductor, antes se descarga (efecto punta). El
punto marca el globo del enunciado.}
\label{fig:sem-globo}
\end{figure}


\begin{observacion}[Interpretación y órdenes de magnitud]\label{obs:sem-globo}
La carga máxima es $Q_{\max}=4\pi\eo R^{2}\Erup=\frac{(5\cdot10^{-2})^{2}\cdot3\cdot10^{6}}{8{,}988\cdot10^{9}}\,\mathrm{C}\simeq0{,}83\,\mu\mathrm{C}$, con densidad
$\sigma_{\max}=\eo\Erup\simeq2{,}7\cdot10^{-5}\,\mathrm{C/m^{2}}$ y presión electrostática $\Pel=\eo\Erup^{2}/2\simeq40\,\mathrm{Pa}$. La relación
$\phi_{\max}=\Erup R$ explica por qué los conductores pequeños o con puntas se descargan antes (\figpanel{fig:sem-globo}{b}): a igual potencial, cuanto menor es
$R$, mayor es el campo en la superficie. Una esfera de $1\,\mathrm{mm}$ no soporta más de $3\,\mathrm{kV}$; por eso los generadores de Van de Graaff tienen
cúpulas grandes y lisas, y los pararrayos, puntas finas (\cref{obs:pararrayos}).
\end{observacion}

% =========================================================================
\subsection{Problema 3.2. Cuatro esferas concéntricas con una batería}\label{sem:cuatro}
% =========================================================================
\begin{problema}[Cuatro superficies esféricas concéntricas]\label{prob:cuatro}
Sean cuatro superficies conductoras esféricas concéntricas $S_1$, $S_2$, $S_3$ y $S_4$ de radios $R_1<R_2<R_3<R_4$. Suponiendo que son inicialmente neutras y
aplicando una diferencia de potencial $V_0$ entre $S_1$ y $S_4$ ($V_0=\phi_1-\phi_4$), calcula los potenciales de las esferas $S_2$ y $S_3$ así como la carga
que adquiere $S_1$.
\end{problema}

\begin{resolucion}
\paso{Idea.} La batería, conectada entre $S_1$ y $S_4$ (a través de pequeños orificios en $S_2$ y $S_3$, cuyo efecto se desprecia), traslada carga de una a otra:
como estaban neutras, $Q_4=-Q_1$. Las esferas $S_2$ y $S_3$ siguen aisladas y neutras. Lo resolvemos de dos formas: con coeficientes de potencial (M3, camino B)
y con la cadena de Gauss (M2).

\paso{Método 1: coeficientes de potencial.} Por la regla \eqref{eq:regla-concentricas}, $p_{ij}=\kc/\max(R_i,R_j)$. Si ponemos una carga $Q$ en la esfera $S_j$,
para la esfera $S_i$ hay dos posibilidades: si está dentro ($R_i<R_j$), está en el interior de la capa cargada y su potencial es $\kc Q/R_j$; si está fuera
($R_i>R_j$), ve la carga como puntual y su potencial es $\kc Q/R_i$. Con $Q_2=Q_3=0$ y $Q_4=-Q_1$:
\[
  \phi_4=p_{41}Q_1+p_{44}Q_4=\frac{\kc}{R_4}\bigl(Q_1+Q_4\bigr)=0 ,
\]
así que la esfera exterior queda al mismo potencial que el infinito (la carga total es nula). La condición de la batería da
\[
  V_0=\phi_1-\phi_4=\bigl(p_{11}-p_{14}\bigr)Q_1=\frac{Q_1}{4\pi\eo}\Bigl(\frac{1}{R_1}-\frac{1}{R_4}\Bigr)=\frac{Q_1}{4\pi\eo}\,\frac{R_4-R_1}{R_1R_4}
  \quad\Longrightarrow\quad\resalta{Q_1=4\pi\eo V_0\,\frac{R_1R_4}{R_4-R_1}.}
\]
Para las esferas intermedias,
\[
  \phi_2=p_{21}Q_1+p_{24}Q_4=\bigl(p_{21}-p_{24}\bigr)Q_1=\frac{Q_1}{4\pi\eo}\Bigl(\frac{1}{R_2}-\frac{1}{R_4}\Bigr)=V_0\,\frac{R_1}{R_2}\,\frac{R_4-R_2}{R_4-R_1},
\]
\[
  \phi_3=\bigl(p_{31}-p_{34}\bigr)Q_1=\frac{Q_1}{4\pi\eo}\Bigl(\frac{1}{R_3}-\frac{1}{R_4}\Bigr)=V_0\,\frac{R_1}{R_3}\,\frac{R_4-R_3}{R_4-R_1}.
\]

\paso{Método 2: cadena de Gauss.} Las esferas neutras $S_2$ y $S_3$ llevan $\mp Q_1$ en sus caras interior y exterior, ambas en el mismo radio, y no modifican el
campo. Por Gauss con esferas concéntricas,
\[
  \vE(\vc{r})=\begin{cases}\vc{0}, & r>R_4\ (\Qenc=Q_1+Q_4=0),\\[3pt] \dfrac{Q_1}{4\pi\eo r^{2}}\,\ur, & R_1<r<R_4,\\[8pt] \vc{0}, & r<R_1 ,\end{cases}
\]
y, como $\phi(r)-\phi(R_4)=\int_r^{R_4}E\,\dif r'$ con $\phi(R_4)=0$,
\[
  \phi(r)=\begin{cases}0, & r\ge R_4,\\[3pt] \dfrac{Q_1}{4\pi\eo}\Bigl(\dfrac1r-\dfrac{1}{R_4}\Bigr), & R_1\le r\le R_4,\\[8pt]
  \dfrac{Q_1}{4\pi\eo}\Bigl(\dfrac{1}{R_1}-\dfrac{1}{R_4}\Bigr)=V_0, & r\le R_1 .\end{cases}
\]
Evaluando en $R_2$ y $R_3$ se recuperan $\phi_2$ y $\phi_3$, y la última línea reproduce $Q_1$.
\end{resolucion}

% =========================================================================
%  Figura (Problema 3.2): cuatro esferas concéntricas con una batería entre
%  S1 y S4. (b) phi(r)/V0 con R_i = i: phi/V0 = (4/r - 1)/3 en [1, 4].
% =========================================================================
\begin{figure}[t]
\centering
\begin{tikzpicture}[font=\small]
  % ---------------------------------------------------------------- (a)
  \begin{scope}[x=1cm, y=1cm]
    \foreach \r/\n in {1.75/4,1.3/3,0.85/2}{\draw[azulprof, line width=0.7pt] (0,0) circle (\r);}
    \shade[ball color=azulvelo] (0,0) circle (0.4);
    \draw[azulprof, line width=0.6pt] (0,0) circle (0.4);
    \foreach \r/\n in {0.4/1,0.85/2,1.3/3,1.75/4}{
      \node[font=\scriptsize, inner sep=0.5pt] at (125:{\r+0.22}) {$S_{\n}$};}
    % hilo por los orificios hasta la batería
    \draw[azulprof, line width=0.6pt] (0.4,0) -- (2.35,0) -- (2.35,-0.35);
    \draw[azulprof, line width=1.1pt] (2.15,-0.35) -- (2.55,-0.35);
    \draw[azulprof, line width=0.6pt] (2.25,-0.47) -- (2.45,-0.47);
    \draw[azulprof, line width=0.6pt] (2.35,-0.47) -- (2.35,-1.0) -- (1.42,-1.0);
    \node[font=\footnotesize, anchor=west] at (2.6,-0.42) {$V_0$};
    \node[font=\scriptsize, anchor=west, text=cargapos] at (2.6,0.0) {$+$};
    \foreach \a in {60,180,240,300}{\node[signopos, font=\tiny\bfseries] at (\a:0.3) {$+$};}
    \foreach \a in {70,200,250}{\node[signoneg, font=\tiny\bfseries] at (\a:1.65) {$-$};}
    \node[font=\footnotesize, text=cargapos] at (0,0.0) {$Q_1$};
    \node[font=\footnotesize, text=carganeg] at (-1.62,-1.62) {$-Q_1$};
    \node at (0,-2.5) {(a)};
  \end{scope}
  % ---------------------------------------------------------------- (b)
  \begin{scope}[xshift=4.4cm, yshift=-2.1cm]
    \begin{axis}[tema, width=8.4cm, height=4.9cm, xmin=0, xmax=5.2, ymin=0, ymax=1.15,
        xlabel={$r$ (con $R_i=i$)}, ylabel={$\phi/V_0$}, xtick={0,1,2,3,4,5},
        xticklabels={$0$,$R_1$,$R_2$,$R_3$,$R_4$,}, ytick={0,0.111,0.333,1}, yticklabels={$0$,$\frac19$,$\frac13$,$1$}]
      \addplot[azulprof, line width=1.2pt] coordinates {(0,1) (1,1)};
      \addplot[azulprof, line width=1.2pt, domain=1:4, samples=60] {(4/x-1)/3};
      \addplot[azulprof, line width=1.2pt] coordinates {(4,0) (5.2,0)};
      \foreach \x in {1,2,3,4}{\edef\tmp{\noexpand\draw[punteado] (axis cs:\x,0) -- (axis cs:\x,1.15);}\tmp}
      \fill[cargapos] (axis cs:2,0.3333) circle (1.8pt) node[above right=-1pt, font=\scriptsize] {$\phi_2$};
      \fill[cargapos] (axis cs:3,0.1111) circle (1.8pt) node[above right=-1pt, font=\scriptsize] {$\phi_3$};
      \fill[cargapos] (axis cs:1,1) circle (1.8pt) node[above right=-1pt, font=\scriptsize] {$\phi_1=V_0$};
      \fill[cargapos] (axis cs:4,0) circle (1.8pt) node[above right=-1pt, font=\scriptsize] {$\phi_4=0$};
    \end{axis}
    \node at (3.4,-1.35) {(b)};
  \end{scope}
\end{tikzpicture}
\caption{Problema 3.2. (a) Cuatro esferas concéntricas; la batería conecta $S_1$ y $S_4$ a través de pequeños orificios en $S_2$ y $S_3$. (b) Potencial con
$R_i=i$: entre $R_1$ y $R_4$ es el de la carga $Q_1$ y fuera de $S_4$ es nulo; las esferas neutras $S_2$ y $S_3$ solo toman el valor del potencial en su radio.}
\label{fig:sem-cuatro}
\end{figure}


\begin{observacion}[Interpretación]\label{obs:sem-cuatro}
Las esferas $S_1$ y $S_4$ están en influencia total: forman un condensador esférico de capacidad $C=Q_1/V_0=4\pi\eo R_1R_4/(R_4-R_1)$, que no depende de
la presencia de $S_2$ y $S_3$. Es un error frecuente pensar que estas modifican el campo: al ser neutras y delgadas, solo adoptan el potencial del radio en el que
están, $V_0>\phi_2>\phi_3>0$ (\figpanel{fig:sem-cuatro}{b}).
\end{observacion}

% =========================================================================
\subsection{Problema 3.3. Condensador esférico: coeficientes}\label{sem:esferico}
% =========================================================================
\begin{problema}[Coeficientes de un condensador esférico]\label{prob:esferico}
Sea un condensador esférico formado por una esfera conductora de radio $a$ rodeada concéntricamente por otra esfera conductora de radio $b$. Calcula los
coeficientes de capacidad e influencia y la capacidad del condensador.
\end{problema}

\begin{resolucion}
\paso{Idea.} Hay dos caminos (método M3). Con conductores concéntricos y cargas fijas, el camino B (coeficientes de potencial y luego inversión) es muy
rápido; el camino A (potenciales impuestos) sirve de comprobación.

\paso{Camino B: coeficientes de potencial.} \emph{Estado $Q_1=1$, $Q_2=0$:} la esfera exterior queda neutra, con $-1$ y $+1$ en sus caras; fuera de la esfera
interior el campo es el de una carga unidad en el centro, luego $\phi_1=\kc/a$ y $\phi_2=\kc/b$: $p_{11}=\kc/a$, $p_{21}=\kc/b$. \emph{Estado $Q_1=0$, $Q_2=1$:} el
campo dentro de la esfera exterior es nulo, y $\phi_1=\phi_2=\kc/b$: $p_{12}=p_{22}=\kc/b$. Así,
\[
  \Pm=\kc\begin{pmatrix}1/a & 1/b\\ 1/b & 1/b\end{pmatrix},\qquad\det\Pm=\kc^{2}\Bigl(\frac{1}{ab}-\frac{1}{b^{2}}\Bigr)=\kc^{2}\,\frac{b-a}{ab^{2}},
\]
\[
  \Cm=\Pm^{-1}=\frac{1}{\det\Pm}\begin{pmatrix}p_{22} & -p_{12}\\ -p_{12} & p_{11}\end{pmatrix}
  =\frac{ab^{2}}{\kc(b-a)}\begin{pmatrix}1/b & -1/b\\ -1/b & 1/a\end{pmatrix}.
\]

\paso{Camino A: potenciales impuestos.} \emph{Estado $\phi_1=1$, $\phi_2=0$:} entre las esferas el campo es el de $Q_1$ en el centro, y
$\phi_1-\phi_2=\kc Q_1\bigl(\frac1a-\frac1b\bigr)=1$, de donde $c_{11}=Q_1=4\pi\eo\frac{ab}{b-a}$. La exterior está a tierra, así que fuera el potencial es nulo y la
carga total también: $c_{21}=Q_2=-Q_1=-c_{11}$. \emph{Estado $\phi_1=0$, $\phi_2=1$:} ahora $\phi_1-\phi_2=\kc Q_1\bigl(\frac1a-\frac1b\bigr)=-1$ da
$c_{12}=Q_1=-4\pi\eo\frac{ab}{b-a}$; y $\phi_2=\kc(Q_1+Q_2)/b=1$ da $Q_1+Q_2=4\pi\eo b$, es decir, $c_{22}=4\pi\eo b-c_{12}$ (\cref{fig:sem-esferico}).
Ambos caminos coinciden:
\[
  \resalta{c_{11}=4\pi\eo\,\frac{ab}{b-a},\qquad c_{12}=c_{21}=-4\pi\eo\,\frac{ab}{b-a},\qquad c_{22}=4\pi\eo\,\frac{b^{2}}{b-a},\qquad C=4\pi\eo\,\frac{ab}{b-a}.}
\]
La capacidad sale de \cref{eq:cap-coef}, $C=1/(p_{11}+p_{22}-2p_{12})=1/[\kc(1/a-1/b)]$, o del \cref{cor:encerrado}, $C=c_{11}$.
\end{resolucion}

% =========================================================================
%  Figura (Problema 3.3): los dos estados del camino A con a = 1, b = 2.
%  phi1 = 1, phi2 = 0: phi = 1 (r <= 1), 2/r - 1 (1..2), 0 (r >= 2)
%  phi1 = 0, phi2 = 1: phi = 0 (r <= 1), 2 - 2/r (1..2), 2/r (r >= 2)
% =========================================================================
\begin{figure}[t]
\centering
\begin{tikzpicture}[font=\small]
  \begin{axis}[tema, name=a, width=7.0cm, height=4.8cm, xmin=0, xmax=4.2, ymin=0, ymax=1.18,
      xlabel={$r/a$}, ylabel={$\phi$}, xtick={0,1,2,3,4}, ytick={0,0.5,1},
      title={\footnotesize $\phi_1=1,\ \phi_2=0$: columna 1 de $\Cm$}]
    \fill[azulvelo!60] (axis cs:0,0) rectangle (axis cs:1,1.18);
    \addplot[azulprof, line width=1.2pt] coordinates {(0,1) (1,1)};
    \addplot[azulprof, line width=1.2pt, domain=1:2, samples=40] {2/x-1};
    \addplot[azulprof, line width=1.2pt] coordinates {(2,0) (4.2,0)};
    \draw[punteado] (axis cs:2,0) -- (axis cs:2,1.18);
    \node[font=\scriptsize, anchor=west, align=left] at (axis cs:2.2,0.65) {$Q_1=c_{11}>0$\\ $Q_2=c_{21}=-c_{11}$};
  \end{axis}
  \node[anchor=north] at ($(a.outer south)+(0,-0.05)$) {(a)};
  \begin{axis}[tema, name=b, at={($(a.east)+(1.4cm,0)$)}, anchor=west, width=7.0cm, height=4.8cm,
      xmin=0, xmax=4.2, ymin=0, ymax=1.18,
      xlabel={$r/a$}, ylabel={$\phi$}, xtick={0,1,2,3,4}, ytick={0,0.5,1},
      title={\footnotesize $\phi_1=0,\ \phi_2=1$: columna 2 de $\Cm$}]
    \fill[azulvelo!60] (axis cs:0,0) rectangle (axis cs:1,1.18);
    \addplot[azulacero, line width=1.2pt] coordinates {(0,0) (1,0)};
    \addplot[azulacero, line width=1.2pt, domain=1:2, samples=40] {2-2/x};
    \addplot[azulacero, line width=1.2pt, domain=2:4.2, samples=40] {2/x};
    \draw[punteado] (axis cs:2,0) -- (axis cs:2,1.18);
    \node[font=\scriptsize, anchor=west, align=left] at (axis cs:2.45,0.38) {$Q_1=c_{12}<0$\\ $Q_2=c_{22}>0$};
  \end{axis}
  \node[anchor=north] at ($(b.outer south)+(0,-0.05)$) {(b)};
\end{tikzpicture}
\caption{Problema 3.3, camino A ($b=2a$). En cada estado se fija el potencial de un conductor a la unidad y el otro a tierra; las cargas que aparecen son una
columna de la matriz de coeficientes de capacidad e influencia. (a) La esfera interior a potencial unidad: todo su flujo acaba en la exterior. (b) La esfera
exterior a potencial unidad: parte de su flujo va al infinito y parte induce carga negativa en la interior, a tierra.}
\label{fig:sem-esferico}
\end{figure}


\begin{observacion}[Comprobaciones obligatorias]\label{obs:sem-esferico}
(i) \emph{Simetría:} $c_{12}=c_{21}$. (ii) \emph{Signos:} $c_{11},c_{22}>0$ y $c_{12}<0$. (iii) \emph{Influencia:} $c_{11}=-c_{21}$, porque la esfera interior está
en influencia total; y $c_{22}+c_{12}=4\pi\eo b>0$, porque parte de las líneas de la exterior van al infinito. (iv) \emph{Límites:} si $b=a+d$ con $d\ll a$,
$C\simeq\eo\,4\pi a^{2}/d$ (condensador plano); si $b\to\infty$, $C\to4\pi\eo a$, la capacidad de la esfera aislada. (v) \emph{Unidades:} $\eo$ por una
longitud, es decir, faradios.
\end{observacion}

% =========================================================================
\subsection{Problema 3.4. Condensador cilíndrico}\label{sem:cilindrico}
% =========================================================================
\begin{problema}[Capacidad por unidad de longitud]\label{prob:cilindrico}
Un condensador cilíndrico se compone de dos cilindros conductores coaxiales de radio interno $a$ y externo $b$. Calcula la capacidad por unidad de longitud
del condensador.
\end{problema}

\begin{resolucion}
\paso{Idea.} Trabajamos con cilindros de longitud $\ell\gg b$ y despreciamos los efectos de borde, de modo que hay simetría cilíndrica. Cargamos las
armaduras con $\pm Q$, es decir, $\pm\lambda=\pm Q/\ell$ por unidad de longitud, y seguimos los cinco pasos del método directo (M2).

\paso{Paso 1: campo entre armaduras.} Por la simetría, $\vE=E(r)\,\urho$. Con una superficie gaussiana cilíndrica coaxial de radio $a<r<b$ y longitud $\ell$, las
tapas no aportan flujo y solo se encierra la carga de la armadura interior:
\[
  2\pi r\ell\,E=\frac{\lambda\ell}{\eo}\quad\Longrightarrow\quad\vE=\frac{\lambda}{2\pi\eo r}\,\urho\qquad(a<r<b),
\]
y el campo es nulo dentro del cilindro interior y fuera del exterior (la carga encerrada es $\lambda-\lambda=0$).

\paso{Paso 2: diferencia de potencial.} Integrando de la armadura positiva a la negativa,
\[
  \DV=\phi_a-\phi_b=-\int_{b}^{a}\vE\cdot\dl=\int_{a}^{b}\frac{\lambda}{2\pi\eo r}\,\dif r=\frac{\lambda}{2\pi\eo}\,\ln\frac{b}{a}.
\]
\paso{Paso 3: capacidad.}
\[
  C=\frac{Q}{\DV}=\frac{\lambda\ell}{\frac{\lambda}{2\pi\eo}\ln(b/a)}\quad\Longrightarrow\quad\resalta{\frac{C}{\ell}=\frac{2\pi\eo}{\ln(b/a)} .}
\]

\paso{Con coeficientes de potencial.} Para un cilindro indefinido no se puede tomar $\phi(\infty)=0$: el potencial de un hilo cargado diverge
logarítmicamente. Tomamos el origen de potenciales a una distancia arbitraria $r_0$ del eje. Un cilindro de radio $R$ con carga $Q$ en una longitud $\ell$
crea entonces (\figpanel{fig:sem-cilindrico}{b})
\[
  \phi(r)=\begin{cases}\dfrac{Q}{2\pi\eo\ell}\,\ln\dfrac{r_0}{r}, & r\ge R,\\[8pt] \dfrac{Q}{2\pi\eo\ell}\,\ln\dfrac{r_0}{R}, & r\le R,\end{cases}
\]
y, razonando como en el caso esférico,
\[
  \Pm=\frac{1}{2\pi\eo\ell}\begin{pmatrix}\ln(r_0/a) & \ln(r_0/b)\\ \ln(r_0/b) & \ln(r_0/b)\end{pmatrix}.
\]
Como $p_{22}=p_{12}$, $p_{11}+p_{22}-2p_{12}=p_{11}-p_{12}$ y
\[
  C=\bigl(p_{11}-p_{12}\bigr)^{-1}=2\pi\eo\ell\Bigl[\ln\frac{r_0}{a}-\ln\frac{r_0}{b}\Bigr]^{-1}=\frac{2\pi\eo\ell}{\ln(b/a)} ,
\]
que no depende de $r_0$, como debe ser. Elegir $r_0=b$ simplifica mucho el problema, porque entonces $\ln(r_0/b)=\ln1=0$.
\end{resolucion}

% =========================================================================
%  Figura (Problema 3.4): condensador cilíndrico.
%  (a) Sección transversal: +lambda en r = a, -lambda en r = b, gaussiana.
%  (b) Potencial de un cilindro cargado respecto de una referencia r0:
%      phi ∝ ln(r0/R) (r <= R), ln(r0/r) (r >= R), con R = 1 y r0 = 3.
% =========================================================================
\begin{figure}[t]
\centering
\begin{tikzpicture}[font=\small]
  % ---------------------------------------------------------------- (a)
  \begin{scope}[x=0.95cm, y=0.95cm]
    \foreach \a in {0,30,...,330}{
      \draw[azulacero, line width=0.55pt, -{Stealth[length=3.5pt]}] (\a:0.78) -- (\a:1.35);
      \draw[azulacero, line width=0.55pt] (\a:1.35) -- (\a:1.85);}
    \draw[azulprof, line width=1.1pt] (0,0) circle (1.95);
    \shade[ball color=azulvelo] (0,0) circle (0.7);
    \draw[azulprof, line width=0.6pt] (0,0) circle (0.7);
    \draw[gauss] (0,0) circle (1.3);
    \foreach \a in {0,45,...,315}{\node[signopos, font=\tiny\bfseries] at (\a:0.55) {$+$};
                                 \node[signoneg, font=\tiny\bfseries] at (\a+22.5:1.82) {$-$};}
    \node[font=\scriptsize, text=cargapos] at (0,0) {$+\lambda$};
    \node[font=\scriptsize, text=carganeg] at (1.9,1.9) {$-\lambda$};
    \draw[-{Stealth[length=3.5pt]}, azulprof, line width=0.5pt] (0,0) -- (-125:0.7);
    \node[font=\scriptsize] at (-0.55,-0.25) {$a$};
    \draw[-{Stealth[length=3.5pt]}, azulprof, line width=0.5pt] (0,0) -- (-55:1.95);
    \node[font=\scriptsize] at (1.0,-1.25) {$b$};
    \node[font=\scriptsize, text=cargapos!85] at (-1.05,1.15) {$r$};
    \node at (0,-2.6) {(a)};
  \end{scope}
  % ---------------------------------------------------------------- (b)
  \begin{scope}[xshift=3.95cm, yshift=-2.15cm]
    \begin{axis}[tema, width=8.0cm, height=4.9cm, xmin=0, xmax=6.2, ymin=-0.75, ymax=1.35,
        axis x line=middle, axis y line=left,
        xlabel={$r$}, ylabel={$\phi\,/\,\frac{Q}{2\pi\eo\ell}$}, xtick={1,3}, xticklabels={$R$,$r_0$}, ytick={-0.5,0,0.5,1},
        xlabel style={anchor=west}]
      \addplot[azulprof, line width=1.2pt] coordinates {(0,1.0986) (1,1.0986)};
      \addplot[azulprof, line width=1.2pt, domain=1:6.2, samples=60] {ln(3/x)};
      \draw[punteado] (axis cs:1,-0.75) -- (axis cs:1,1.35);
      \fill[cargapos] (axis cs:3,0) circle (1.8pt);
      \node[font=\scriptsize, anchor=south west, text=grispizarra] at (axis cs:3.05,0.05) {referencia $\phi(r_0)=0$};
      \node[font=\scriptsize, anchor=west, text=grispizarra] at (axis cs:1.15,1.12) {$\ln(r_0/R)$};
    \end{axis}
    \node at (3.2,-1.35) {(b)};
  \end{scope}
\end{tikzpicture}
\caption{Problema 3.4. (a) Sección transversal del condensador cilíndrico (longitud $\ell$): la gaussiana coaxial de radio $a<r<b$ solo encierra la carga de la
armadura interior. (b) Potencial de un cilindro cargado de radio $R$: como diverge en el infinito, se refiere a una distancia arbitraria $r_0$; las diferencias
de potencial, y por tanto la capacidad, no dependen de esa elección.}
\label{fig:sem-cilindrico}
\end{figure}


\begin{observacion}[Interpretación]\label{obs:sem-cilindrico}
Con placas próximas, $b=a+d$ y $d\ll a$, $\ln(1+d/a)\simeq d/a$ y $C/\ell\simeq2\pi\eo a/d=\eo(A/\ell)/d$, con $A/\ell=2\pi a$: de nuevo el condensador plano. Los
coeficientes de potencial respecto de $r_0$ ya no cumplen todas las propiedades del \cref{teo:pij-prop} (si $r_0<b$, $p_{22}<0$), porque estas usaban
$\phi(\infty)=0$; las diferencias de potencial y la capacidad, en cambio, son independientes de la referencia.
\end{observacion}

% =========================================================================
\subsection{Problema 3.5. Placas planas paralelas}\label{sem:plano}
% =========================================================================
\begin{problema}[Condensador plano y reparto de carga por caras]\label{prob:plano}
Sean dos conductores planos paralelos de superficie $A$ y separados una distancia $d$. Despreciando el efecto de bordes, calcula la capacidad del
condensador. Si la carga de cada conductor es $Q_1$ y $Q_2$, calcula las cargas en cada una de las caras de los conductores.
\end{problema}

\begin{resolucion}
\paso{Idea.} Si las dimensiones de las placas son grandes frente a $d$, cada cara se comporta como un plano indefinido cargado, que crea un campo
$\sigma/2\eo$ a cada lado, dirigido hacia fuera del plano si $\sigma>0$ (método M4). Con dos placas hay cuatro caras, con cargas $q_1,q_2,q_3,q_4$ de
izquierda a derecha (\cref{fig:sem-plano}).

\paso{Paso 1: condensador ($Q_1=Q$, $Q_2=-Q$).} Las cargas $\pm Q$ se sitúan en las caras enfrentadas, con densidades $\pm\sigma=\pm Q/A$. El campo entre las
placas puede obtenerse de dos formas. \emph{Por superposición:} entre las placas los campos de ambas se suman, $\sigma/2\eo+\sigma/2\eo=\sigma/\eo$, y fuera se
anulan. \emph{Por Gauss}, con una caja que tiene una tapa dentro del metal ($\vE=\vc{0}$) y otra entre las placas:
\[
  E\,A_{\mathrm{caja}}=\frac{\sigma A_{\mathrm{caja}}}{\eo}\quad\Longrightarrow\quad E=\frac{\sigma}{\eo}=\frac{Q}{\eo A}.
\]
El campo es uniforme, de modo que $\DV=E\,d=\frac{Qd}{\eo A}$ y
\[
  \resalta{C=\frac{Q}{\DV}=\eo\,\frac{A}{d} .}
\]

\paso{Paso 2: cargas cualesquiera $Q_1$ y $Q_2$.} Tomamos como positivo el sentido hacia la derecha. Un punto interior a la placa~1 está a la derecha de la cara
$q_1$ y a la izquierda de $q_2$, $q_3$ y $q_4$; uno interior a la placa~2, a la derecha de $q_1,q_2,q_3$ y a la izquierda de $q_4$. Imponiendo $\vE=\vc{0}$ en
ambos y la conservación de la carga en cada placa:
\[
  \begin{aligned}
    q_1-q_2-q_3-q_4&=0 &&\text{(campo nulo en la placa 1)},\\
    q_1+q_2+q_3-q_4&=0 &&\text{(campo nulo en la placa 2)},\\
    q_1+q_2&=Q_1,\qquad q_3+q_4=Q_2 &&\text{(conservación de la carga)}.
  \end{aligned}
\]
Restando las dos primeras, $2(q_2+q_3)=0$, es decir, $q_3=-q_2$; sumándolas, $2(q_1-q_4)=0$, es decir, $q_1=q_4$. Con la conservación,
$q_1+q_2=Q_1$ y $q_1-q_2=Q_2$, de donde
\[
  \resalta{q_1=q_4=\frac{Q_1+Q_2}{2}\ \ \text{(caras exteriores)},\qquad q_2=-q_3=\frac{Q_1-Q_2}{2}\ \ \text{(caras interiores)}.}
\]
\end{resolucion}

% =========================================================================
%  Figura (Problema 3.5): placas planas paralelas.
%  (a) Condensador ideal: +sigma y -sigma en las caras enfrentadas, campo
%      sigma/eps0 entre las placas y nulo fuera; caja de Gauss.
%  (b) Cargas cualesquiera: cuatro caras q1..q4 y campo en cada región.
% =========================================================================
\begin{figure}[t]
\centering
\begin{tikzpicture}[font=\small]
  % ---------------------------------------------------------------- (a)
  \begin{scope}
    \fill[top color=azulvelo!80!white, bottom color=metal, draw=azulprof, line width=0.6pt] (-1.75,-1.7) rectangle (-1.4,1.7);
    \fill[top color=azulvelo!80!white, bottom color=metal, draw=azulprof, line width=0.6pt] (1.4,-1.7) rectangle (1.75,1.7);
    \foreach \y in {-1.5,-1.1,...,1.6}{\node[signopos, font=\tiny\bfseries] at (-1.28,\y) {$+$};
                                       \node[signoneg, font=\tiny\bfseries] at (1.28,\y) {$-$};}
    \foreach \y in {-1.35,-0.85,0.85,1.35}{\draw[lcampo] (-1.1,\y) -- (1.1,\y);}
    \node[font=\footnotesize, fill=white, inner sep=1.5pt] at (0.25,-0.2) {$E=\dfrac{\sigma}{\eo}$};
    \draw[gaussazul] (-1.62,0.15) rectangle (-0.75,0.55);
    \node[font=\scriptsize, text=azulacero, anchor=west] at (-0.7,0.35) {gaussiana};
    \node[font=\footnotesize] at (-2.35,-0.9) {$E=0$};
    \node[font=\footnotesize] at (2.35,-0.9) {$E=0$};
    \node[font=\footnotesize, text=cargapos] at (-1.58,-2.0) {$+Q$};
    \node[font=\footnotesize, text=carganeg] at (1.58,-2.0) {$-Q$};
    \draw[{Stealth[length=3pt]}-{Stealth[length=3pt]}, grispizarra] (-1.4,1.95) -- (1.4,1.95) node[midway, above, font=\footnotesize] {$d$};
    \node at (0,-3.15) {(a)};
  \end{scope}
  % ---------------------------------------------------------------- (b)
  \begin{scope}[xshift=8.6cm]
    \fill[top color=azulvelo!80!white, bottom color=metal, draw=azulprof, line width=0.6pt] (-2.2,-1.7) rectangle (-1.6,1.7);
    \fill[top color=azulvelo!80!white, bottom color=metal, draw=azulprof, line width=0.6pt] (1.6,-1.7) rectangle (2.2,1.7);
    \foreach \x/\q in {-2.2/q_1,-1.6/q_2,1.6/q_3,2.2/q_4}{
      \draw[cargapos, line width=1.3pt] (\x,-1.7) -- (\x,1.7);
      \node[font=\footnotesize, anchor=south] at (\x,1.72) {$\q$};}
    \node[font=\scriptsize, rotate=90] at (-1.9,0) {$\vE=\vc{0}$};
    \node[font=\scriptsize, rotate=90] at (1.9,0) {$\vE=\vc{0}$};
    \node[font=\footnotesize] at (-1.9,-2.0) {$Q_1$};
    \node[font=\footnotesize] at (1.9,-2.0) {$Q_2$};
    \draw[-{Stealth[length=4pt]}, azulacero, line width=0.8pt] (-2.45,1.0) -- (-3.25,1.0);
    \node[font=\scriptsize, anchor=north] at (-2.85,0.85) {$\dfrac{Q_1+Q_2}{2\eo A}$};
    \draw[-{Stealth[length=4pt]}, azulacero, line width=0.8pt] (-0.4,1.0) -- (0.4,1.0);
    \node[font=\scriptsize, anchor=north] at (0,0.85) {$\dfrac{Q_1-Q_2}{2\eo A}$};
    \draw[-{Stealth[length=4pt]}, azulacero, line width=0.8pt] (2.45,1.0) -- (3.25,1.0);
    \node[font=\scriptsize, anchor=north] at (2.85,0.85) {$\dfrac{Q_1+Q_2}{2\eo A}$};
    \node[font=\scriptsize, align=center] at (0,-2.45) {$q_1=q_4=\tfrac12(Q_1+Q_2)$\\[3pt] $q_2=-q_3=\tfrac12(Q_1-Q_2)$};
    \node at (0,-3.15) {(b)};
  \end{scope}
\end{tikzpicture}
\caption{Problema 3.5. (a) Condensador plano ideal: la carga está en las caras enfrentadas y el campo, uniforme, solo existe entre las placas; la caja de Gauss
con una tapa dentro del metal da $E=\sigma/\eo$. (b) Placas con cargas cualesquiera: las caras exteriores se reparten la carga total y las interiores llevan
cargas opuestas; se indica el campo en cada región (para $Q_1>Q_2$ y $Q_1+Q_2>0$).}
\label{fig:sem-plano}
\end{figure}


\begin{observacion}[Interpretación]\label{obs:sem-plano}
Las caras exteriores comparten la carga total a partes iguales, y las enfrentadas llevan cargas opuestas: solo estas forman el condensador. El campo es
$\frac{Q_1-Q_2}{2\eo A}$ entre las placas y $\frac{Q_1+Q_2}{2\eo A}$ fuera, hacia fuera a cada lado (\cref{tab:sem-plano}). Con $Q_1=Q$ y $Q_2=-Q$ se recupera el
condensador ideal: caras exteriores descargadas y campo nulo fuera. Si, por ejemplo, $Q_2=0$, cada cara de la placa neutra lleva $\mp Q_1/2$: la placa neutra
«parte» el campo pero no lo anula.
\end{observacion}

\begin{table}[ht]
\centering
\caption{Campo de cada placa del condensador plano ideal ($\pm\sigma$) en las tres regiones (positivo hacia la derecha).}
\label{tab:sem-plano}
\begin{tabular}{@{}lccc@{}}
  \toprule
  & Izquierda & Entre placas & Derecha\\
  \midrule
  Placa $+\sigma$ & $-\dfrac{\sigma}{2\eo}$ & $+\dfrac{\sigma}{2\eo}$ & $+\dfrac{\sigma}{2\eo}$\\[8pt]
  Placa $-\sigma$ & $+\dfrac{\sigma}{2\eo}$ & $+\dfrac{\sigma}{2\eo}$ & $-\dfrac{\sigma}{2\eo}$\\[6pt]
  \midrule
  Total & $0$ & $\dfrac{\sigma}{\eo}$ & $0$\\[6pt]
  \bottomrule
\end{tabular}
\end{table}

% =========================================================================
\subsection{Problema 3.6. Tres esferas alineadas}\label{sem:alineadas}
% =========================================================================
\begin{problema}[Tres esferas alineadas, laterales a tierra]\label{prob:alineadas}
Tres esferas conductoras de radio $R$ se sitúan alineadas con sus centros separados entre sí $d$ ($d\gg R$). Las laterales se conectan a tierra y a la central
se le suministra una carga $Q$. Halla el potencial de esta última.
\end{problema}

\begin{resolucion}
\paso{Idea.} Como $d\gg R$, aplicamos el método M5: las esferas se ven entre sí como cargas puntuales y los coeficientes de potencial son los de la regla
\eqref{eq:regla-lejanas}. Numeramos 1 (izquierda), 2 (central) y 3 (derecha), con $d_{12}=d_{23}=d$ y $d_{13}=2d$.

\paso{Paso 1: matriz de coeficientes de potencial.}
\[
  \Pm=\kc\begin{pmatrix}1/R & 1/d & 1/(2d)\\ 1/d & 1/R & 1/d\\ 1/(2d) & 1/d & 1/R\end{pmatrix}.
\]
\paso{Paso 2: clasificación.} Las esferas 1 y 3 tienen como dato $\phi=0$ y como incógnita su carga; la 2 tiene como dato $Q_2=Q$ y como incógnita $\phi_2$. Por
simetría, $Q_1=Q_3\equiv q$.

\paso{Paso 3: esferas a tierra.}
\[
  0=\phi_1=\kc\Bigl(\frac{q}{R}+\frac{Q}{d}+\frac{q}{2d}\Bigr)\quad\Longrightarrow\quad q\,\frac{2d+R}{2dR}=-\frac{Q}{d}
  \quad\Longrightarrow\quad q=-\frac{2QR}{2d+R}.
\]
\paso{Paso 4: potencial buscado.}
\[
  \phi_2=\kc\Bigl(\frac{Q}{R}+\frac{2q}{d}\Bigr)=\frac{Q}{4\pi\eo}\Bigl[\frac1R-\frac{4R}{d(2d+R)}\Bigr]
  \quad\Longrightarrow\quad\resalta{\phi_2=\frac{Q}{4\pi\eo}\,\frac{2d^{2}+Rd-4R^{2}}{Rd\,(2d+R)} .}
\]
\end{resolucion}

% =========================================================================
%  Figura (Problema 3.6): tres esferas alineadas, laterales a tierra.
%  (b) phi2/(kQ/R) frente a d/R: 1 - 4/(x(2x+1)) y la aproximación 1 - 2/x^2.
% =========================================================================
\begin{figure}[t]
\centering
\begin{tikzpicture}[font=\small]
  % ---------------------------------------------------------------- (a)
  \begin{scope}
    \foreach \x/\n in {-2.2/1,0/2,2.2/3}{
      \shade[ball color=azulvelo] (\x,0) circle (0.38);
      \draw[azulprof, line width=0.6pt] (\x,0) circle (0.38);
      \node[font=\footnotesize] at (\x,0) {$\n$};}
    \foreach \x in {-2.2,2.2}{\draw[azulprof, line width=0.6pt] (\x,-0.38) -- (\x,-0.8); \pic at (\x,-0.8) {tierra};
                              \node[font=\footnotesize, text=carganeg] at (\x,0.65) {$q$};}
    \node[font=\footnotesize, text=cargapos] at (0,0.65) {$Q$};
    \draw[{Stealth[length=3pt]}-{Stealth[length=3pt]}, grispizarra] (-2.2,-1.45) -- (0,-1.45) node[midway, below, font=\footnotesize] {$d$};
    \draw[{Stealth[length=3pt]}-{Stealth[length=3pt]}, grispizarra] (0,-1.45) -- (2.2,-1.45) node[midway, below, font=\footnotesize] {$d$};
    \draw[-{Stealth[length=3pt]}, azulprof, line width=0.5pt] (0,0) ++(-130:0.02) -- ++(-130:0.36);
    \node[font=\scriptsize] at (-0.5,-0.55) {$R$};
    \node[font=\footnotesize, text=grispizarra] at (0,1.2) {$\phi_1=\phi_3=0$};
    \node at (0,-2.45) {(a)};
  \end{scope}
  % ---------------------------------------------------------------- (b)
  \begin{scope}[xshift=4.0cm, yshift=-2.0cm]
    \begin{axis}[tema, width=7.6cm, height=4.6cm, xmin=3, xmax=20, ymin=0.8, ymax=1.03,
        xlabel={$d/R$}, ylabel={$\phi_2\,/\,(\kc Q/R)$}, ytick={0.8,0.85,0.9,0.95,1},
        legend style={at={(0.97,0.05)}, anchor=south east}]
      \addplot[grismedio, line width=0.8pt, densely dashed, domain=3:20] {1};
      \addlegendentry{esfera aislada}
      \addplot[azulprof, line width=1.2pt, domain=3:20, samples=80] {1-4/(x*(2*x+1))};
      \addlegendentry{resultado}
      \addplot[azulcielo, line width=1.1pt, densely dotted, domain=3:20, samples=80] {1-2/x^2};
      \addlegendentry{$1-2R^{2}/d^{2}$}
    \end{axis}
    \node at (2.95,-1.35) {(b)};
  \end{scope}
\end{tikzpicture}
\caption{Problema 3.6. (a) Esfera central con carga $Q$ entre dos esferas a tierra. (b) Potencial de la esfera central normalizado al de una esfera aislada: las
cargas negativas inducidas en las laterales lo rebajan en una cantidad de orden $(R/d)^{2}$.}
\label{fig:sem-alineadas}
\end{figure}


\begin{observacion}[Interpretación y orden de la aproximación]\label{obs:sem-alineadas}
Las esferas a tierra adquieren carga negativa, $q\simeq-QR/d$, que rebaja el potencial de la central respecto del de una esfera aislada, $\kc Q/R$:
\[
  \phi_2=\frac{Q}{4\pi\eo R}\Bigl[1-\frac{4R^{2}}{d(2d+R)}\Bigr]\simeq\frac{Q}{4\pi\eo R}\Bigl(1-\frac{2R^{2}}{d^{2}}\Bigr).
\]
La corrección es pequeña, de orden $(R/d)^{2}$ (\figpanel{fig:sem-alineadas}{b}). El resultado solo tiene sentido al orden más bajo en $R/d$: los términos de orden
superior (como el $R^{3}/d^{3}$ que aparece al desarrollar la expresión exacta) no son significativos dentro de la aproximación de esferas lejanas, porque se
han despreciado efectos del mismo orden (la deformación de la carga de cada esfera por las vecinas).
\end{observacion}

% =========================================================================
\subsection{Problema 3.7. Tres esferas en triángulo, a tierra por etapas}\label{sem:triangulo}
% =========================================================================
\begin{problema}[Tres esferas a tierra una a una]\label{prob:triangulo}
Sean tres esferas conductoras de radio $a$ situadas en los vértices de un triángulo equilátero de lado $b$ ($b\gg a$). Las tres esferas tienen carga $q$.
Sucesivamente, una a una se llevan a potencial de tierra mediante una conexión con un hilo conductor muy delgado hasta que se alcanza el equilibrio
electrostático, desconectándolas posteriormente. Calcula las cargas $q_1$, $q_2$ y $q_3$ de las esferas al final del proceso.
\end{problema}

\begin{resolucion}
\paso{Idea.} Método M5 por etapas (\cref{fig:sem-triangulo}). En cada etapa, solo la esfera conectada a tierra tiene su carga como incógnita; las demás están
aisladas y conservan la carga que tenían, incluidas las ya desconectadas: \emph{lo que se congela es la carga, no el potencial}.

\paso{Paso 1: matriz de coeficientes de potencial.} Como $b\gg a$, cada esfera se ve desde las otras como una carga puntual. El coeficiente propio es el
potencial de una esfera aislada con carga unidad, $p_{ii}\simeq\kc/a$, y el mutuo, el que crea una carga unidad a distancia $b$, $p_{ij}\simeq\kc/b$; por la
simetría del triángulo, los tres mutuos son iguales:
\[
  \Pm=\kc\begin{pmatrix}1/a & 1/b & 1/b\\ 1/b & 1/a & 1/b\\ 1/b & 1/b & 1/a\end{pmatrix},\qquad
  \phi_i=\kc\Bigl(\frac{Q_i}{a}+\frac{Q_j+Q_l}{b}\Bigr),
\]
con $j$ y $l$ las otras dos esferas.

\paso{Paso 2: clasificación, etapa por etapa.} En la etapa 1 la esfera 1 está a tierra y $Q_2=Q_3=q$; en la etapa 2, la esfera 2 está a tierra, $Q_1=q_1$
(congelada) y $Q_3=q$; en la etapa 3, la esfera 3 está a tierra y $Q_1=q_1$, $Q_2=q_2$ (congeladas).

\paso{Paso 3: una ecuación por etapa.} Solo se usa la ecuación de la esfera a tierra:
\[
  \text{Etapa 1:}\quad0=\phi_1=\kc\Bigl(\frac{q_1}{a}+\frac{q+q}{b}\Bigr)\quad\Longrightarrow\quad q_1=-\frac{2qa}{b},
\]
\[
  \text{Etapa 2:}\quad0=\phi_2=\kc\Bigl(\frac{q_2}{a}+\frac{q_1+q}{b}\Bigr)\quad\Longrightarrow\quad
  q_2=-\frac{a}{b}\bigl(q+q_1\bigr)=-\frac{qa}{b}\Bigl(1-\frac{2a}{b}\Bigr)=-\frac{qa\,(b-2a)}{b^{2}},
\]
\[
  \text{Etapa 3:}\quad0=\phi_3=\kc\Bigl(\frac{q_3}{a}+\frac{q_1+q_2}{b}\Bigr)\quad\Longrightarrow\quad
  q_3=-\frac{a}{b}\bigl(q_1+q_2\bigr)=\frac{a}{b}\,\frac{qa}{b^{2}}\bigl[2b+(b-2a)\bigr]=\frac{qa^{2}(3b-2a)}{b^{3}} .
\]
En resumen,
\[
  \resalta{q_1=-\frac{2qa}{b},\qquad q_2=-\frac{qa\,(b-2a)}{b^{2}},\qquad q_3=\frac{qa^{2}(3b-2a)}{b^{3}} .}
\]
\end{resolucion}

% =========================================================================
%  Figura (Problema 3.7): proceso por etapas y cargas en cada etapa para
%  a/b = 0,1 (q = 1): inicial (1, 1, 1); etapa 1 (-0.2, 1, 1);
%  etapa 2 (-0.2, -0.08, 1); etapa 3 (-0.2, -0.08, 0.028).
% =========================================================================
\newcommand{\triEtapa}[5]{% #1 título, #2 #3 #4 rótulos de carga, #5 esfera a tierra (0 = ninguna)
  \coordinate (v1) at (-0.8,0.45); \coordinate (v2) at (0.8,0.45); \coordinate (v3) at (0,-0.94);
  \draw[grismedio, line width=0.4pt, densely dashed] (v1) -- (v2) -- (v3) -- cycle;
  \foreach \v/\n in {v1/1,v2/2,v3/3}{
    \ifnum\n=#5 \shade[ball color=azulcielo] (\v) circle (0.2); \else \shade[ball color=azulvelo] (\v) circle (0.2); \fi
    \draw[azulprof, line width=0.5pt] (\v) circle (0.2);
    \node[font=\tiny] at (\v) {\n};}
  \node[font=\scriptsize, anchor=south] at ($(v1)+(0,0.22)$) {#2};
  \node[font=\scriptsize, anchor=south] at ($(v2)+(0,0.22)$) {#3};
  \node[font=\scriptsize, anchor=west] at ($(v3)+(0.24,0)$) {#4};
  \ifnum#5=1 \draw[azulprof, line width=0.5pt] (v1) ++(-0.2,0) -- ++(-0.3,0) -- ++(0,-0.25); \pic at ($(v1)+(-0.5,-0.25)$) {tierra};\fi
  \ifnum#5=2 \draw[azulprof, line width=0.5pt] (v2) ++(0.2,0) -- ++(0.3,0) -- ++(0,-0.25); \pic at ($(v2)+(0.5,-0.25)$) {tierra};\fi
  \ifnum#5=3 \draw[azulprof, line width=0.5pt] (v3) ++(-0.2,0) -- ++(-0.3,0) -- ++(0,-0.25); \pic at ($(v3)+(-0.5,-0.25)$) {tierra};\fi
  \node[font=\footnotesize, anchor=north] at (0,-1.45) {#1};}

\begin{figure}[t]
\centering
\begin{tikzpicture}[font=\small]
  \begin{scope}[xshift=0cm]\triEtapa{Inicial}{$q$}{$q$}{$q$}{0}\end{scope}
  \begin{scope}[xshift=3.6cm]\triEtapa{Etapa 1: $\phi_1=0$}{$q_1$}{$q$}{$q$}{1}\end{scope}
  \begin{scope}[xshift=7.2cm]\triEtapa{Etapa 2: $\phi_2=0$}{$q_1$}{$q_2$}{$q$}{2}\end{scope}
  \begin{scope}[xshift=10.8cm]\triEtapa{Etapa 3: $\phi_3=0$}{$q_1$}{$q_2$}{$q_3$}{3}\end{scope}
  \node[font=\scriptsize, text=grispizarra] at (5.4,1.45) {lado $b$, radio $a$};
  \node at (5.4,-2.15) {(a)};
  \begin{scope}[yshift=-6.4cm, xshift=1.7cm]
    \begin{axis}[tema, ybar, bar width=7pt, width=11.8cm, height=4.6cm, ymin=-0.3, ymax=1.15,
        symbolic x coords={Inicial,Etapa 1,Etapa 2,Etapa 3}, xtick=data, ylabel={$Q_i/q$},
        enlarge x limits=0.15, axis x line*=bottom, axis y line*=left, ytick={-0.2,0,0.5,1},
        legend style={at={(0.5,1.03)}, anchor=south, legend columns=3}]
      \addplot[fill=azulprof, draw=azulprof] coordinates {(Inicial,1) (Etapa 1,-0.2) (Etapa 2,-0.2) (Etapa 3,-0.2)};
      \addlegendentry{esfera 1}
      \addplot[fill=azulacero, draw=azulacero] coordinates {(Inicial,1) (Etapa 1,1) (Etapa 2,-0.08) (Etapa 3,-0.08)};
      \addlegendentry{esfera 2}
      \addplot[fill=azulcielo, draw=azulcielo] coordinates {(Inicial,1) (Etapa 1,1) (Etapa 2,1) (Etapa 3,0.028)};
      \addlegendentry{esfera 3}
    \end{axis}
    \node at (5.2,-1.2) {(b)};
  \end{scope}
\end{tikzpicture}
\caption{Problema 3.7. (a) Esquema del proceso: en cada etapa se conecta a tierra una esfera (sombreada) y después se desconecta; las demás conservan su carga.
(b) Cargas de las tres esferas tras cada etapa para $a/b=0{,}1$: cada esfera conectada a tierra pierde casi toda su carga, y la última acaba con una pequeña
carga positiva porque sus vecinas ya son negativas.}
\label{fig:sem-triangulo}
\end{figure}


\begin{observacion}[Comprobaciones e interpretación]\label{obs:sem-triangulo}
\emph{Patrón común:} la esfera que se conecta a tierra adquiere $-a/b$ veces la suma de las cargas actuales de las otras dos, $Q_i^{\mathrm{final}}=-\frac{a}{b}(Q_j+Q_l)$.
\emph{Signos:} $q_1<0$ porque las otras dos son positivas; $q_2<0$ porque $q+q_1=q(1-2a/b)>0$ al ser $b\gg a$; y $q_3>0$ porque las otras dos son ya negativas.
\emph{Tamaños:} $\abs{q_1},\abs{q_2}\sim qa/b$ y $q_3\sim qa^{2}/b^{2}$: cada esfera conectada a tierra pierde casi toda su carga, y lo que le queda es lo que
inducen las vecinas lejanas. \emph{Lo que se congela es la carga:} al terminar la etapa 2, la esfera 1 ya no está a potencial cero,
\[
  \phi_1=\kc\Bigl(\frac{q_1}{a}+\frac{q_2+q}{b}\Bigr)=\kc\,\frac{q_2-q}{b}\neq0 :
\]
al conectar la esfera 2 su carga cambia y eso modifica el potencial de la 1, pero como la 1 está aislada su carga no puede variar. Por último, los términos como
$2a/b$ dentro de $q_2$ y $q_3$ provienen de propagar la aproximación de una etapa a la siguiente, y coinciden con la solución de la relación.
\end{observacion}

% =========================================================================
\subsection{Ejemplo modelo: tres esferas concéntricas, la interior a tierra}\label{sem:modelo}
% =========================================================================
\begin{problema}[Método de los coeficientes de potencial]\label{prob:modelo}
Tres esferas conductoras concéntricas de radios $a$, $b=2a$ y $c=3a$ tienen cargas $q_a=q$, $q_b=2q$ y $q_c=3q$. ¿Cuánto varía el potencial de la esfera
externa cuando la interna se conecta a tierra?
\end{problema}

\begin{resolucion}
\paso{Paso 1: matriz.} Con la regla \eqref{eq:regla-concentricas},
\[
  \Pm=\frac{\kc}{a}\begin{pmatrix}1 & 1/2 & 1/3\\ 1/2 & 1/2 & 1/3\\ 1/3 & 1/3 & 1/3\end{pmatrix}.
\]
\paso{Paso 2: estado inicial.} Todas aisladas, con cargas conocidas:
$\phi_c=p_{ca}q_a+p_{cb}q_b+p_{cc}q_c=\frac{\kc}{3a}(q+2q+3q)=\frac{2\kc q}{a}$.

\paso{Paso 3: clasificación tras conectar.} Para $a$, el dato es $\phi_a=0$ y la incógnita $q_a'$; las esferas $b$ y $c$ siguen aisladas, con $q_b=2q$ y
$q_c=3q$.

\paso{Paso 4: ecuación de la esfera a tierra.}
\[
  0=\phi_a=p_{aa}q_a'+p_{ab}q_b+p_{ac}q_c=\kc\Bigl(\frac{q_a'}{a}+\frac{2q}{2a}+\frac{3q}{3a}\Bigr)\quad\Longrightarrow\quad q_a'=-2q .
\]
Por el hilo pasa a tierra una carga $q_a-q_a'=3q$.

\paso{Paso 5: variación del potencial.} Como $\Pm$ solo depende de la geometría, si en un proceso cambian algunas cargas,
$\Delta\phi_i=\sum_jp_{ij}\,\Delta Q_j$, y solo hay que sumar sobre los conductores cuya carga cambia:
\[
  \Delta\phi_c=p_{ca}\bigl(q_a'-q_a\bigr)=\frac{\kc}{3a}(-2q-q)\quad\Longrightarrow\quad\resalta{\Delta\phi_c=-\frac{q}{4\pi\eo a}.}
\]
\end{resolucion}

\begin{observacion}[Comprobación con Gauss]\label{obs:sem-modelo}
Fuera de todo, la carga total pasa de $6q$ a $3q$, así que $\phi_c=\kc Q_{\mathrm{tot}}/c$ pasa de $\kc\,6q/3a=2\kc q/a$ a $\kc\,3q/3a=\kc q/a$: coincide. También se
comprueba $\phi_a'=0$ con la cadena de Gauss (\figpanel{fig:sem-modelo}{b}):
$\phi_a'=\phi_c'+\kc q_a'\bigl(\frac1a-\frac1b\bigr)+\kc(q_a'+q_b)\bigl(\frac1b-\frac1c\bigr)=\frac{\kc q}{a}-\frac{2\kc q}{a}\cdot\frac12+0=0$.
\emph{Trampa habitual:} al conectar a tierra cambia la carga del conductor conectado, no la de los aislados; pero cambian \emph{todos} los potenciales.
\end{observacion}

% =========================================================================
%  Figura (ejemplo modelo): tres esferas concéntricas a, 2a, 3a con cargas
%  q, 2q, 3q. Potencial (en unidades de kq/a) antes y después de conectar la
%  interior a tierra:
%   antes:   3 (r<=1), 2 + 1/r (1..2), 1 + 3/r (2..3), 6/r (r>=3)
%   después: 0 (r<=1), 2 - 2/r (1..2), 1 (2..3), 3/r (r>=3)
% =========================================================================
\begin{figure}[t]
\centering
\begin{tikzpicture}[font=\small]
  % ---------------------------------------------------------------- (a)
  \begin{scope}[x=0.62cm, y=0.62cm]
    \draw[azulprof, line width=0.8pt] (0,0) circle (3);
    \draw[azulprof, line width=0.8pt] (0,0) circle (2);
    \shade[ball color=azulvelo] (0,0) circle (1);
    \draw[azulprof, line width=0.6pt] (0,0) circle (1);
    \node[font=\scriptsize] at (0,0.25) {$q$};
    \node[font=\scriptsize] at (0,-0.3) {$a$};
    \node[font=\scriptsize, fill=white, inner sep=0.5pt] at (90:2) {$2q$};
    \node[font=\scriptsize, fill=white, inner sep=0.5pt] at (90:3) {$3q$};
    \draw[azulprof, line width=0.5pt] (1,0) -- (3.6,0) -- (3.6,-0.5);
    \pic at (3.6,-0.5) {tierra};
    \node[fondo, font=\scriptsize, anchor=south, text=grispizarra] at (2.3,0.05) {hilo};
    \node at (0,-4.1) {(a)};
  \end{scope}
  % ---------------------------------------------------------------- (b)
  \begin{scope}[xshift=4.4cm, yshift=-2.25cm]
    \begin{axis}[tema, width=8.4cm, height=5.0cm, xmin=0, xmax=5, ymin=0, ymax=3.3,
        xlabel={$r/a$}, ylabel={$\phi\,/\,(\kc q/a)$}, xtick={0,1,2,3,4,5}, ytick={0,1,2,3},
        legend style={at={(0.99,0.97)}, anchor=north east}]
      \addplot[azulprof, line width=1.2pt] coordinates {(0,3) (1,3)};
      \addlegendentry{antes}
      \addplot[azulprof, line width=1.2pt, domain=1:2, samples=30, forget plot] {2+1/x};
      \addplot[azulprof, line width=1.2pt, domain=2:3, samples=30, forget plot] {1+3/x};
      \addplot[azulprof, line width=1.2pt, domain=3:5, samples=30, forget plot] {6/x};
      \addplot[azulcielo, line width=1.2pt, densely dashed] coordinates {(0,0) (1,0)};
      \addlegendentry{interior a tierra}
      \addplot[azulcielo, line width=1.2pt, densely dashed, domain=1:2, samples=30, forget plot] {2-2/x};
      \addplot[azulcielo, line width=1.2pt, densely dashed, forget plot] coordinates {(2,1) (3,1)};
      \addplot[azulcielo, line width=1.2pt, densely dashed, domain=3:5, samples=30, forget plot] {3/x};
      \foreach \x in {1,2,3}{\edef\tmp{\noexpand\draw[punteado] (axis cs:\x,0) -- (axis cs:\x,3.3);}\tmp}
      \draw[{Stealth[length=3pt]}-{Stealth[length=3pt]}, cargapos] (axis cs:3.12,1) -- (axis cs:3.12,2)
        node[midway, right, font=\scriptsize] {$\Delta\phi_c$};
    \end{axis}
    \node at (3.6,-1.4) {(b)};
  \end{scope}
\end{tikzpicture}
\caption{Ejemplo modelo. (a) Tres esferas concéntricas de radios $a$, $2a$ y $3a$; la interior se conecta a tierra. (b) Potencial antes y después de la conexión:
la esfera interior pasa de $q$ a $-2q$, el potencial de la externa baja de $2\kc q/a$ a $\kc q/a$ y, entre $2a$ y $3a$, el campo se anula porque la carga
encerrada es $-2q+2q=0$.}
\label{fig:sem-modelo}
\end{figure}


% -------------------------------------------------------------------------
\subsection{Resumen de resultados y lista de control}\label{sem:resumen}
% -------------------------------------------------------------------------
\begin{table}[ht]
\centering
\caption{Resultados de los seminarios del Tema 3.}
\label{tab:resultados}
\small
\begin{tabular}{@{}clc@{}}
  \toprule
  Problema & Resultado & Método\\
  \midrule
  3.1 & $\phi_{\max}=\Erup R=150\,\mathrm{kV}$ & M1\\[3pt]
  3.2 & $Q_1=4\pi\eo V_0\dfrac{R_1R_4}{R_4-R_1}$,\quad $\phi_{2,3}=\dfrac{Q_1}{4\pi\eo}\Bigl(\dfrac{1}{R_{2,3}}-\dfrac{1}{R_4}\Bigr)$ & M2, M3\\[10pt]
  3.3 & $c_{11}=-c_{12}=4\pi\eo\dfrac{ab}{b-a}$,\quad $c_{22}=4\pi\eo\dfrac{b^{2}}{b-a}$,\quad $C=4\pi\eo\dfrac{ab}{b-a}$ & M3\\[10pt]
  3.4 & $C/\ell=2\pi\eo/\ln(b/a)$ & M2\\[3pt]
  3.5 & $C=\eo A/d$;\quad caras: $q_1=q_4=\tfrac12(Q_1+Q_2)$,\ $q_2=-q_3=\tfrac12(Q_1-Q_2)$ & M4\\[4pt]
  3.6 & $\phi_2=\dfrac{Q}{4\pi\eo}\,\dfrac{2d^{2}+Rd-4R^{2}}{Rd\,(2d+R)}$ & M5\\[10pt]
  3.7 & $q_1=-\dfrac{2qa}{b}$,\quad $q_2=-\dfrac{qa(b-2a)}{b^{2}}$,\quad $q_3=\dfrac{qa^{2}(3b-2a)}{b^{3}}$ & M5\\[10pt]
  Modelo & $\Delta\phi_c=-q/(4\pi\eo a)$ & M3\\
  \bottomrule
\end{tabular}
\end{table}

\needspace{20\baselineskip}
\begin{receta}[Lista de control]
\begin{itemize}
  \item ¿He asignado la carga de cada cara con la regla de las caras?
  \item ¿He distinguido \emph{a tierra} ($\phi=0$, pero $Q\neq0$ en general) de \emph{descargado} ($Q=0$, pero $\phi\neq0$ en general)?
  \item En los $p_{ij}$, ¿los demás conductores están descargados? En los $c_{ij}$, ¿están a tierra?
  \item ¿Son simétricas las matrices? ¿Tienen los signos correctos ($p_{ij}>0$, $c_{ii}>0$, $c_{ij}<0$)?
  \item ¿Se cumple $c_{jj}\ge-\sum_{i\neq j}c_{ij}$, con igualdad en influencia total?
  \item ¿Coincide la capacidad calculada directamente con $1/(p_{11}+p_{22}-2p_{12})$?
  \item En procesos por etapas, ¿he congelado la carga de los conductores ya desconectados?
  \item Un hilo fino solo iguala potenciales: ¿he despreciado su carga y su influencia?
\end{itemize}
\end{receta}
